% !PRINTED 416
\weberpage{416}
If $p' > 1$, and hence $\pi > 1$, then the term with the
$(\nu-1)$th power of the unknown, whose coefficient is equal to the negative sum of the roots, is missing from this equation. We therefore
have the theorem:

\gesperrt{The sum of all primitive $n$th roots of unity,
when $n$ is a higher power of a prime,
is always equal to zero.}

For the general formation of $X_n$ we shall apply a recurrent
procedure; we assume $X_n$ as already formed,
denote by $p$ a prime which does not divide $n$, by $p'$,
as above, the $(\pi-1)$th power of $p$, and now form $X_{np'p}$,
in which of course $p'$ may also be equal to 1.

If we denote by $r$ the \gesperrt{primitive} $n$th roots of unity,
by $\alpha$ every root of unity of degree $pp'$, and by $\alpha'$ every root of unity of degree $p'$, then one obtains all the primitive roots of unity of degree $npp'$ by taking away from all the
$r\alpha$ the $r\alpha'$ (\S.\,132, II). But the $r\alpha$ are
the roots of the equation:
\[
X_n(x^{pp'}) = 0,
\]
because the quantities
\[
(r\alpha)^{pp'} = r^{pp'},
\]
agree, apart from their order, with the $r$ themselves
(\S.\,132, IV). Likewise the $r\alpha'$ are the roots of the equation
\[
X_n(x^{p'}) = 0;
\]
and from this it follows that
\weq{9}{X_{npp'}(x) = \frac{X_n(x^{pp'})}{X_n(x^{p'})}.}
The value $n=1$ forms no exception to this,
if we understand by $X_1$ the function $x-1$. According to this
we can form $X_n$ comparatively simply in all cases.
We consider some special cases.

Let us suppose that $n$ contains only two distinct primes $p$, $q$, so that
\[
n = pp'qq',
\]
then the formulae (8) and (9) yield:
\[
X_n = \frac{(x^n-1)\left(x^{\frac{n}{pq}}-1\right)}{\left(x^{\frac{n}{p}}-1\right)\left(x^{\frac{n}{q}}-1\right)},
\]
%CONT>>
